We estimate enrollment-weighted average effects of Medicaid coverage on utilization and spending on GLP-1 medications using a stacked difference-in-differences strategy that is robust to treatment effect heterogeneity in staggered adoption settings \citep{wing2024stacked, goodman2021difference, dechaisemartin2020two}. Our main outcome variables are measures of Medicaid prescriptions and spending per 100 enrollee-months for obesity-indicated and diabetes/cardiovascular-indicated GLP-1 medications. We also examine consumer spending at companies known to sell a mix of compounded and branded GLP-1 medications at low prices, which may be an important source in the absence of insurance coverage.

We use $s$ and $t$ to index states and calendar quarters. $A_s$ is the quarter $s$ starts covering at least one obesity-indicated GLP-1 medication; never-adopting states have $A_s = NT$. $Y_{st}(a)$ and $Y_{st}(0)$ are potential outcomes for the same state and quarter with and without coverage, and $Y_{st} = Y_{st}(0) +\sum_a \bp{Y_{st}(a)-Y_{st}(0)} \times \one{A_s = a}$ is the observed outcome. Let $\Omega_{\kappa}$ represent the set of adoption events that occur early enough to be followed for and event window with $\kappa_{pre}$ pre-periods and $\kappa_{post}$ post-periods. Our \emph{target estimand} is an average of enrollment-weighted group-time effects across a fixed set of adoption groups over the event window:
%
\begin{equation*}
  \theta_{\Omega_{\kappa}}^{e} = \sum_{a \in \Omega_{\kappa}} \operatorname{ATT}(a,\, a+e) \times \frac{M_a}{M_{\Omega_{\kappa}}}
\end{equation*}
%
The aggregate effect is an enrollment weighted average of group-time effects. For adoption group $a$, $\operatorname{ATT}(a,\, a+e) = \E[M_s^{\kappa}]{Y_{s,a+e}(a) - Y_{s,a+e}(0)}{A_s = a}$ is the enrollment-weighted average treatment effect of adopting coverage in period $a$ on outcomes experienced in period $a+e$.\footnote{The expectation is taken over the cross-state distribution of total enrollee-months over the event window. For state $s$,  $M_s^{\kappa} = \sum_{t=a-\kappa_{pre}}^{a+\kappa_{post}} M_{st}$, where $M_{st}$ is the state's total enrollee-months of Medicaid coverage in the quarter.} Each adoption group's effect is weighted by the the group's share of total enrollee-months in the event window, $\frac{M_a}{M_{\Omega_{\kappa}}}$.\footnote{Specifically, let $M_a = \sum_{s|A_s=a} M_s^{\kappa}$ be total enrollee-months across all states in adoption group $a$. Then $M_{\Omega_{\kappa}} = \sum_{a\in\Omega_{\kappa}} M_a$ is total enrollee-months in the study population. The weights sum to one across the adoption groups.} Since the study population of states is fixed across event time periods, changes in $\theta_{\Omega_{\kappa}}^{e}$ across event times reflect treatment effect dynamics rather than compositional shifts. 

We estimate $\theta_{\Omega_{\kappa}}^{e}$ using stacked difference-in-differences with never-treated states serving as clean controls. To construct the analytic sample, we build a sub-experimental dataset for each adoption cohort. The dataset of sub-experiment $a$ consists of the treated states with $A_s = a$, the never-treated control states, and observations from time periods inside the event window. We vertically concatenate the sub-experimental datasets into a single stacked dataset in which $Y_{sae}$ is the observed outcome for state $s$ in sub-experiment $a$ at event time $e$. We estimate the aggregate effects using a weighted least squares regression that is saturated in event time and treatment status:
%
\begin{equation*}
  Y_{sae} = \alpha_0 + \alpha_1 \one{A_s = a} + \sum_{h \ne -1} \gamma_h \one{e = h} + \sum_{h \ne -1} \theta_h \bp{\one{A_s = a} \times \one{e = h}} + \varepsilon_{sae},
\end{equation*}
%
The regressions are weighted to correct for differences in the share of treated and control observations across the sub-experiments and to weight states according to enrollment.\footnote{The weights are derived in \citet{wing2024stacked}. In our case, let $N_a^{treat}$ and $N_a^{control}$ represent the number of treated and control states in sub-experiment $a$. Then $N_{\Omega_{\kappa}}^{treat}$ and $N_{\Omega_{\kappa}}^{control}$ are the total number of treated and control states across the study population. The corrective weights are $Q_{sa} = \frac{\sfrac{M_a}{M_{\Omega_{\kappa}}}}{\sfrac{N_a^{\text{treat}}}{N_{\Omega_{\kappa}}^{\text{treat}}}}$ if $A_s = a$, and $Q_{sa} = \frac{\sfrac{M_a}{M_{\Omega_{\kappa}}}}{\sfrac{N_a^{\text{control}}}{N_{\Omega_{\kappa}}^{\text{control}}}}$ if $A_s = NT$.}  Standard errors are clustered at the state level. The stacked difference-in-differences coefficients identify $\theta_{\Omega_{\kappa}}^{e}$ under three assumptions within each sub-experiment: (i) \emph{common trends}, (ii) \emph{no-anticipation}; and (iii) \emph{no spillovers}. The event study coefficients from pre-treatment event times should equal zero if the common trend and no-anticipation assumption hold, providing a partial test of these restrictions. 

There is a tradeoff between the length of the event window and the number of states that can be included in the analysis: for states that expanded coverage recently, not enough time has elapsed for a long follow up. We consider two designs that resolve the tradeoff in different ways. The \emph{10-state design} allows for 4 pre-treatment quarters and 7 post-treatment quarters and includes all states that covered at least one GLP-1 obesity drug by 2023-Q4, The \emph{7-state design} restricts the sample to states that covered a GLP-1 obesity drug by 2023-Q1, allowing for a 4 pre-treatment quarters and 10 post-treatment quarters. Table~\ref{tab:design} lists the treated states in the two designs. In both designs, the 33 states that had not adopted coverage by the end of 2025-Q2 serve as controls. 
